% ============================================================================
%  ESTADÍSTICA INFERENCIAL — Sesión 2
%  Semana 2: 17 al 22 de agosto de 2026   (lunes 17: festivo)
%  1.2 Reglas de conteo, combinaciones y permutaciones
%  Compilar: pdflatex Est-Inferencial-S02.tex  (dos veces)
% ============================================================================
\documentclass[aspectratio=169,10pt]{beamer}

\usepackage[utf8]{inputenc}
\IfFileExists{spanish.ldf}{\usepackage[spanish,es-noshorthands]{babel}}{}
\usepackage{amsmath,amssymb}
\usepackage{booktabs}
\usepackage{tikz}
\usepackage{listings}
\usetikzlibrary{positioning}

\usetheme{Uniminuto}

\uniasignatura{Estadística inferencial}
\uniprofesor{Andrés Felipe Puerta Velez}
\unicorreo{andres.puerta-v@uniminuto.edu.co}
\uninrc{CBASBL061}
\unisemana{Sesión 2 — Semana 2 · 17 al 22 de agosto de 2026}

\newcommand{\listatiny}{\setbeamerfont{itemize/enumerate body}{size=\tiny}\tiny}

\newcommand{\fx}[1]{{\ttfamily\color{uniAzul}#1}}

% ------------------------------------------------------------- CÓDIGO R ----
\lstdefinestyle{Rcode}{
  language=R,
  basicstyle=\ttfamily\scriptsize,
  keywordstyle=\color{uniAzul}\bfseries,
  commentstyle=\color{uniGris}\itshape,
  stringstyle=\color{uniTitulo},
  showstringspaces=false, numbers=none, breaklines=true,
  columns=fullflexible,
  backgroundcolor=\color{uniCremaClaro},
  frame=leftline, framerule=1.2pt, rulecolor=\color{uniNaranja},
  xleftmargin=0.35cm, framexleftmargin=0.35cm,
  framexrightmargin=0.15cm, framextopmargin=2pt, framexbottommargin=2pt,
  aboveskip=0.15cm, belowskip=0.15cm
}
\lstdefinestyle{Rmini}{
  language=R,
  basicstyle=\ttfamily\tiny,
  keywordstyle=\color{uniAzul}\bfseries,
  commentstyle=\color{uniGris}\itshape,
  stringstyle=\color{uniTitulo},
  showstringspaces=false, numbers=none, breaklines=true,
  columns=fullflexible,
  backgroundcolor=\color{uniCremaClaro},
  frame=leftline, framerule=1.2pt, rulecolor=\color{uniNaranja},
  xleftmargin=0.35cm, framexleftmargin=0.35cm,
  framexrightmargin=0.15cm, framextopmargin=2pt, framexbottommargin=2pt,
  aboveskip=0.12cm, belowskip=0.12cm
}
\lstdefinestyle{Rout}{
  basicstyle=\ttfamily\tiny\color{black},
  showstringspaces=false, numbers=none, breaklines=true,
  columns=fullflexible,
  backgroundcolor=\color{uniGris!10},
  frame=leftline, framerule=1.2pt, rulecolor=\color{uniGris},
  xleftmargin=0.35cm, framexleftmargin=0.35cm,
  framexrightmargin=0.15cm, framextopmargin=2pt, framexbottommargin=2pt,
  aboveskip=0.12cm, belowskip=0.12cm
}

\begin{document}

\portada

% ============================================================== ORDEN DEL DÍA
\begin{frame}{Orden del día}
  \begin{enumerate}
    \item Llamado a lista
    \item Repaso: dónde quedamos en la sesión anterior
    \item \azul{Tema 1.2: reglas de conteo, permutaciones y combinaciones}
    \item Principio de la multiplicación y regla de la suma
    \item Factorial, permutaciones y combinaciones
    \item Ejercicios en clase
    \item Cierre y ejercicios propuestos
  \end{enumerate}
\end{frame}

% ================================================================== REPASO ==
\begin{frame}{Repaso: dónde quedamos}
  En la sesión 1 llegamos a que, si todos los resultados son igualmente
  probables,
  \[ P(A)=\frac{\#A}{\#\Omega} \]

  \vspace{0.2cm}
  \begin{alertblock}{El problema de hoy}
    \footnotesize
    Con un dado podíamos escribir $\Omega=\{1,\dots,6\}$ y contar a mano. Pero
    si el experimento es ``asignar un código de inventario de 2 letras y 3
    dígitos'', enumerar todos los resultados es imposible: hay cientos de
    miles. \destacar{Necesitamos contar sin enumerar.}
  \end{alertblock}

  \vspace{0.2cm}
  {\footnotesize\color{uniGris} De eso se ocupan las reglas de conteo: son las
  herramientas para calcular $\#A$ y $\#\Omega$ cuando los conjuntos son
  grandes\,[2].}
\end{frame}

% ============================================================== UNIDAD 1 ====
\seccion{1.2 Reglas de conteo}

% ------------------------------------------------------- MULTIPLICACIÓN ----
\begin{frame}{Principio de la multiplicación}
  \begin{block}{Enunciado}
    \small
    Si una tarea se realiza en $k$ etapas, y la etapa $i$ se puede hacer de
    $n_i$ maneras (sin importar lo elegido antes), entonces la tarea completa
    se puede hacer de
    \[ n_1 \times n_2 \times \cdots \times n_k \quad\text{maneras.} \]
  \end{block}

  \vspace{0.15cm}
  \begin{columns}[T,onlytextwidth]
    \begin{column}{0.46\textwidth}
      \footnotesize
      \textbf{Ejemplo.} Un menú ofrece 3 entradas, 4 platos fuertes y 2
      postres. Una comida completa se arma en 3 etapas:
      \[ 3\times 4\times 2 = 24 \text{ comidas distintas.} \]
    \end{column}
    \begin{column}{0.50\textwidth}
      \centering
      \begin{tikzpicture}[scale=0.63,every node/.style={font=\scriptsize},
                          level distance=13mm,sibling distance=7mm]
        \node[circle,draw=uniAzul,fill=uniAzul!12,inner sep=1.4pt] (r) {}
          child {node[circle,draw=uniNaranja,fill=uniNaranja!25,inner sep=1.4pt]{}
            child {node[circle,fill=uniAzul,inner sep=1pt]{}}
            child {node[circle,fill=uniAzul,inner sep=1pt]{}}}
          child {node[circle,draw=uniNaranja,fill=uniNaranja!25,inner sep=1.4pt]{}
            child {node[circle,fill=uniAzul,inner sep=1pt]{}}
            child {node[circle,fill=uniAzul,inner sep=1pt]{}}}
          child {node[circle,draw=uniNaranja,fill=uniNaranja!25,inner sep=1.4pt]{}
            child {node[circle,fill=uniAzul,inner sep=1pt]{}}
            child {node[circle,fill=uniAzul,inner sep=1pt]{}}};
        \node[anchor=west,text=uniGris] at (2.6,0) {etapa 1: 3 opciones};
        \node[anchor=west,text=uniGris] at (2.6,-1.3) {etapa 2: 2 opciones};
        \node[anchor=west,text=uniGris] at (2.6,-2.6) {$3\times 2=6$ caminos};
      \end{tikzpicture}
    \end{column}
  \end{columns}
\end{frame}

% ---------------------------------------------------------- REGLA SUMA -----
\begin{frame}{Regla de la suma: cuándo se suma y cuándo se multiplica}
  \begin{columns}[T,onlytextwidth]
    \begin{column}{0.48\textwidth}
      \begin{block}{Se \emph{multiplica}}
        \footnotesize
        Cuando las elecciones son \azul{etapas sucesivas} de una misma tarea:
        se hace lo uno \textbf{y} lo otro.

        \vspace{0.15cm}
        ``Elijo entrada \textbf{y} plato fuerte'' $\rightarrow 3\times 4$
      \end{block}
    \end{column}
    \begin{column}{0.48\textwidth}
      \begin{exampleblock}{Se \emph{suma}}
        \footnotesize
        Cuando las opciones son \azul{alternativas excluyentes}: se hace lo uno
        \textbf{o} lo otro, nunca ambos.

        \vspace{0.15cm}
        ``Voy en bus (4 rutas) \textbf{o} en metro (2 líneas)''
        $\rightarrow 4+2=6$
      \end{exampleblock}
    \end{column}
  \end{columns}

  \vspace{0.35cm}
  \begin{alertblock}{Conexión con la sesión anterior}
    \footnotesize
    La regla de la suma solo vale si las alternativas son \destacar{mutuamente
    excluyentes}. Si no lo son, hay que descontar la intersección, exactamente
    como en $\#(A\cup B)=\#A+\#B-\#(A\cap B)$.
  \end{alertblock}
\end{frame}

% ------------------------------------------------------------ EJERCICIO 1 --
\begin{frame}{Ejercicio 1 — en clase}
  \begin{alertblock}{Individual (6 minutos)}
    Una empresa asigna códigos de inventario formados por \destacar{2 letras}
    (alfabeto de 26) seguidas de \destacar{3 dígitos} (del 0 al 9).

    \vspace{0.3cm}
    \begin{enumerate}
      \item ¿Cuántos códigos distintos existen si se permite repetir letras y
            dígitos?
      \item ¿Cuántos si las dos letras deben ser diferentes?
      \item ¿Cuántos códigos del punto 1 terminan en dígito par?
    \end{enumerate}
  \end{alertblock}
\end{frame}

\begin{frame}{Ejercicio 1 — solución}
  \small
  \begin{block}{Se aplica el principio de la multiplicación en cada caso}
    \begin{align*}
      \text{1.}\quad & 26\times 26\times 10\times 10\times 10 = 26^{2}\cdot 10^{3}
                       = \azul{676\,000}\\[0.25em]
      \text{2.}\quad & 26\times 25\times 10^{3} = \azul{650\,000}\\[0.25em]
      \text{3.}\quad & 26^{2}\times 10\times 10\times 5 = \azul{338\,000}
    \end{align*}
  \end{block}

  \vspace{0.15cm}
  \begin{exampleblock}{Observaciones}
    \footnotesize
    En el punto 2 la segunda letra tiene solo 25 opciones porque no puede
    repetir la primera: eso es una \azul{permutación}. En el punto 3 la última
    etapa tiene 5 opciones $\{0,2,4,6,8\}$; note que $338\,000$ es exactamente
    la mitad de $676\,000$, como era de esperarse.
  \end{exampleblock}
\end{frame}

% -------------------------------------------------------------- FACTORIAL --
\begin{frame}{Factorial y permutaciones}
  \begin{block}{Factorial}
    \small
    \[ n! = n\times(n-1)\times\cdots\times 2\times 1,
       \qquad\text{con } 0!=1 \text{ por convención.} \]
  \end{block}

  \vspace{0.15cm}
  \begin{block}{Permutación de los $n$ objetos}
    \small
    El número de formas de \azul{ordenar} $n$ objetos distintos en fila es $n!$.

    \vspace{0.1cm}
    \footnotesize
    \emph{Por qué:} para el primer puesto hay $n$ opciones, para el segundo
    $n-1$, y así sucesivamente. Por el principio de la multiplicación,
    $n(n-1)\cdots 1 = n!$.
  \end{block}

  \vspace{0.15cm}
  \begin{exampleblock}{Ejemplo}
    \footnotesize
    Ocho empleados se ordenan para una foto: $8! = 40\,320$ formas.
    Seis personas en una fila: $6! = 720$.
  \end{exampleblock}
\end{frame}

% --------------------------------------------------------- PERMUTACIONES ---
\begin{frame}{Permutaciones de $r$ objetos entre $n$}
  \begin{block}{Fórmula}
    \[ P(n,r)=\frac{n!}{(n-r)!}
       = \underbrace{n\,(n-1)\cdots(n-r+1)}_{r\ \text{factores}} \]
    Se usa cuando se eligen $r$ de $n$ objetos \destacar{y el orden importa}.
  \end{block}

  \vspace{0.2cm}
  \begin{columns}[T,onlytextwidth]
    \begin{column}{0.48\textwidth}
      \begin{exampleblock}{Ejemplo}
        \footnotesize
        De 8 empleados se eligen presidente, secretario y tesorero:
        \[ P(8,3)=\frac{8!}{5!}=8\cdot 7\cdot 6 = 336 \]
      \end{exampleblock}
    \end{column}
    \begin{column}{0.48\textwidth}
      \begin{alertblock}{Con repetición}
        \footnotesize
        Si los objetos se pueden repetir (como los dígitos de un código), el
        conteo es $n^{r}$, no $P(n,r)$.
        \[ 10^{3}=1000 \text{ códigos de 3 dígitos} \]
      \end{alertblock}
    \end{column}
  \end{columns}
\end{frame}

% ------------------------------------------------ PERMUTACIONES REPETIDAS --
\begin{frame}{Permutaciones con objetos repetidos}
  \begin{block}{Fórmula}
    \small
    Si entre $n$ objetos hay $n_1$ iguales de un tipo, $n_2$ de otro, \dots,
    $n_k$ del último, el número de ordenamientos \emph{distinguibles} es
    \[ \frac{n!}{n_1!\;n_2!\cdots n_k!} \]
  \end{block}

  \vspace{0.2cm}
  \begin{exampleblock}{Ejemplo: la palabra BANANA}
    \footnotesize
    Tiene 6 letras, con 3 letras A, 2 letras N y 1 letra B:
    \[ \frac{6!}{3!\,2!\,1!}=\frac{720}{12}=60 \text{ palabras distintas.} \]
    Se divide porque intercambiar dos letras A \azul{no} produce una palabra
    nueva.
  \end{exampleblock}
\end{frame}

% ------------------------------------------------------------ EJERCICIO 2 --
\begin{frame}{Ejercicio 2 — en clase}
  \begin{alertblock}{En parejas (7 minutos)}
    Un área de una empresa tiene \destacar{8 empleados}.

    \vspace{0.3cm}
    \begin{enumerate}
      \item ¿De cuántas formas se puede elegir un comité con presidente,
            secretario y tesorero?
      \item ¿De cuántas formas se pueden elegir 3 empleados para un comité
            \emph{sin} cargos definidos?
      \item ¿De cuántas formas se pueden ordenar los 8 para una fotografía?
      \item ¿Por qué la respuesta del punto 1 es mayor que la del punto 2?
    \end{enumerate}
  \end{alertblock}
\end{frame}

\begin{frame}{Ejercicio 2 — solución}
  \begin{columns}[T,onlytextwidth]
    \begin{column}{0.50\textwidth}
      \begin{block}{Cálculos}
        \small
        \begin{align*}
          \text{1.}\;& P(8,3)=8\cdot 7\cdot 6 = \azul{336}\\[0.25em]
          \text{2.}\;& C(8,3)=\frac{8!}{3!\,5!} = \azul{56}\\[0.25em]
          \text{3.}\;& 8! = \azul{40\,320}
        \end{align*}
      \end{block}
    \end{column}
    \begin{column}{0.46\textwidth}
      \begin{exampleblock}{Punto 4}
        \footnotesize
        Porque en el punto 1 el \azul{orden importa}: los mismos tres nombres
        producen $3!=6$ comités distintos según quién sea presidente.

        \vspace{0.2cm}
        \[ 336 = 56 \times 3! \]
        Esa es exactamente la relación entre permutaciones y combinaciones.
      \end{exampleblock}
    \end{column}
  \end{columns}
\end{frame}

% -------------------------------------------------------- COMBINACIONES ----
\begin{frame}{Combinaciones}
  \begin{block}{Fórmula}
    \[ C(n,r)=\binom{n}{r}=\frac{n!}{r!\,(n-r)!}
       = \frac{P(n,r)}{r!} \]
    Se usa cuando se eligen $r$ de $n$ objetos \destacar{y el orden no importa}.
  \end{block}

  \vspace{0.2cm}
  \begin{columns}[T,onlytextwidth]
    \begin{column}{0.48\textwidth}
      \begin{exampleblock}{Ejemplo}
        \footnotesize
        De 12 jugadores se arma un equipo de 5:
        \[ \binom{12}{5}=792 \]
      \end{exampleblock}
    \end{column}
    \begin{column}{0.48\textwidth}
      \begin{alertblock}{Propiedad útil}
        \footnotesize
        \[ \binom{n}{r}=\binom{n}{n-r} \]
        Elegir 4 de 10 para incluir es lo mismo que elegir 6 de 10 para
        excluir: $\binom{10}{4}=\binom{10}{6}=210$.
      \end{alertblock}
    \end{column}
  \end{columns}
\end{frame}

% -------------------------------------------------------- CÓMO DECIDIR -----
\begin{frame}{¿Permutación o combinación?}
  \begin{block}{La única pregunta que hay que hacerse}
    \centering\small
    \destacar{¿Cambiar el orden produce un resultado diferente?}\\[0.2em]
    Si la respuesta es \emph{sí} $\rightarrow$ permutación. \quad
    Si es \emph{no} $\rightarrow$ combinación.
  \end{block}

  \vspace{0.2cm}
  \footnotesize
  \begin{center}
  \begin{tabular}{@{}lll@{}}
    \toprule
    \textbf{Situación} & \textbf{¿Importa el orden?} & \textbf{Se usa} \\
    \midrule
    Cargos en un comité (presidente, secretario)  & Sí & $P(n,r)$ \\
    Integrantes de un comité sin cargos           & No & $C(n,r)$ \\
    Podio de una competencia (1.º, 2.º, 3.º)      & Sí & $P(n,r)$ \\
    Números ganadores de una lotería              & No & $C(n,r)$ \\
    Clave de acceso de 4 dígitos                  & Sí, y con repetición & $n^{r}$ \\
    Muestra de 3 productos de un lote             & No & $C(n,r)$ \\
    \bottomrule
  \end{tabular}
  \end{center}
\end{frame}

% ------------------------------------------------------------- RESUMEN -----
\begin{frame}{Resumen: las cuatro fórmulas de conteo}
  \footnotesize
  \begin{center}
  \begin{tabular}{@{}llll@{}}
    \toprule
    \textbf{Situación} & \textbf{¿Orden?} & \textbf{¿Repetición?} &
    \textbf{Fórmula} \\
    \midrule
    Etapas sucesivas          & ---   & --- & $n_1\cdot n_2\cdots n_k$ \\
    Ordenar los $n$ objetos   & Sí    & No  & $n!$ \\
    Elegir $r$ de $n$         & Sí    & No  & $P(n,r)=\dfrac{n!}{(n-r)!}$ \\[0.35em]
    Elegir $r$ de $n$         & Sí    & Sí  & $n^{r}$ \\[0.15em]
    Elegir $r$ de $n$         & No    & No  & $C(n,r)=\dfrac{n!}{r!\,(n-r)!}$ \\[0.35em]
    Ordenar con repetidos     & Sí    & ---  & $\dfrac{n!}{n_1!\cdots n_k!}$ \\
    \bottomrule
  \end{tabular}
  \end{center}

  \vspace{0.15cm}
  {\scriptsize\color{uniGris} En calculadora científica: $P(n,r)$ suele
  aparecer como \texttt{nPr} y $C(n,r)$ como \texttt{nCr}\,[3].}
\end{frame}

% ------------------------------------------------------------ EJERCICIO 3 --
\begin{frame}{Ejercicio 3 — en clase}
  \begin{alertblock}{Aplicado (10 minutos)}
    En una caja hay \destacar{12 productos}, de los cuales \destacar{4 son
    defectuosos}. Se extraen 3 al azar, de una sola vez (el orden no importa).

    \vspace{0.25cm}
    \begin{enumerate}
      \item ¿De cuántas formas se pueden extraer 3 productos?
      \item ¿En cuántas de esas formas los 3 resultan buenos?
      \item ¿Cuál es la probabilidad de que los 3 sean buenos?
      \item ¿Cuál es la probabilidad de que haya exactamente 1 defectuoso?
    \end{enumerate}
  \end{alertblock}
\end{frame}

\begin{frame}{Ejercicio 3 — solución}
  \begin{columns}[T,onlytextwidth]
    \begin{column}{0.50\textwidth}
      \begin{block}{Conteos}
        \small
        \begin{align*}
          \#\Omega &= \binom{12}{3} = \azul{220}\\[0.3em]
          \#A &= \binom{8}{3} = \azul{56}
        \end{align*}
        \footnotesize
        (los 8 productos buenos, tomados de a 3)
      \end{block}
    \end{column}
    \begin{column}{0.46\textwidth}
      \begin{exampleblock}{Probabilidades}
        \footnotesize
        \[ P(\text{3 buenos})=\frac{56}{220}=0{,}2545 \]
        \[ P(\text{exactamente 1 defectuoso})
           =\frac{\binom{4}{1}\binom{8}{2}}{\binom{12}{3}}
           =\frac{4\cdot 28}{220}=\frac{112}{220}=0{,}5091 \]
      \end{exampleblock}
    \end{column}
  \end{columns}

  \vspace{0.1cm}
  {\footnotesize\color{uniGris} El punto 4 combina las dos reglas: se
  \emph{multiplica} para armar cada caso favorable (1 defectuoso \textbf{y} 2
  buenos) y se divide por el total de casos posibles.}
\end{frame}

% ------------------------------------------------------------- PUENTE ------
\begin{frame}{Del conteo a la probabilidad}
  Las reglas de conteo son el motor de la regla de Laplace: sirven para
  calcular el numerador y el denominador de
  \[ P(A)=\frac{\#A}{\#\Omega} \]

  \vspace{0.2cm}
  \begin{exampleblock}{Ejemplo completo}
    \footnotesize
    Se lanza una moneda 5 veces. El espacio muestral tiene $2^{5}=32$
    secuencias. Las que tienen exactamente 3 caras son
    $\binom{5}{3}=10$, porque basta elegir \emph{en cuáles} de los 5
    lanzamientos salió cara. Entonces
    \[ P(\text{exactamente 3 caras})=\frac{10}{32}=0{,}3125 \]
  \end{exampleblock}

  \vspace{0.1cm}
  {\footnotesize\color{uniGris} En la sesión 3 formalizaremos la regla de
  Laplace y los enfoques para asignar probabilidades.}
\end{frame}

% ------------------------------------------------------------- CIERRE ------
\begin{frame}[fragile]{Conteo, en R}
  \scriptsize
\begin{lstlisting}[style=Rmini]
factorial(5)             # 5!
choose(10, 3)            # combinaciones C(10,3)
factorial(10)/factorial(7)   # permutaciones P(10,3)
prod(4, 3, 2)            # principio multiplicativo

combn(c("A","B","C","D"), 2)   # las combinaciones, una por columna
nrow(expand.grid(1:6, 1:6))    # espacio muestral de dos dados
\end{lstlisting}
\begin{lstlisting}[style=Rout]
[1] 120
[1] 120
[1] 720
[1] 24
     [,1] [,2] [,3] [,4] [,5] [,6]
[1,] "A"  "A"  "A"  "B"  "B"  "C"
[2,] "B"  "C"  "D"  "C"  "D"  "D"
[1] 36
\end{lstlisting}

  \vspace{-0.1cm}
  \begin{alertblock}{Para qué sirve de verdad}
    \tiny
    \fx{choose()} evita el error más común de la sesión: equivocarse al
    simplificar factoriales. Y \fx{combn()} \destacar{muestra} las
    combinaciones, no solo las cuenta: es la mejor forma de convencerse de que
    el orden no importa.
  \end{alertblock}
\end{frame}

\begin{frame}{Cierre de la sesión}
  \begin{block}{Lo que debe quedar claro hoy}
    \begin{itemize}
      \item Etapas sucesivas se multiplican; alternativas excluyentes se suman.
      \item Si el orden importa: $n!$, $P(n,r)$ o $n^{r}$ según haya o no
            repetición.
      \item Si el orden no importa: $C(n,r)$, y siempre $C(n,r)=P(n,r)/r!$.
      \item El conteo alimenta a $P(A)=\#A/\#\Omega$.
    \end{itemize}
  \end{block}

  \vspace{0.25cm}
  \begin{alertblock}{Trabajo independiente para la próxima sesión}
    \footnotesize
    Resolver los \destacar{20 ejercicios de la lámina siguiente}; las
    respuestas están en la lámina posterior, para que usted mismo se corrija.
    Practicar las teclas \texttt{nPr} y \texttt{nCr} de la calculadora y
    revisar la referencia [3].
  \end{alertblock}
\end{frame}

% -------------------------------------------------- EJERCICIOS PROPUESTOS --
\begin{frame}{Ejercicios propuestos}
  \vspace{0.1cm}
  \listatiny
  \vspace{0.15cm}
  \begin{columns}[T,onlytextwidth]
    \begin{column}{0.47\textwidth}
      \begin{enumerate}\setlength{\itemsep}{0.16cm}
        \item $5!$
        \item $8!/5!$
        \item $P(7,2)$
        \item $C(7,2)$
        \item $C(10,4)$
        \item $C(10,6)$
        \item Placas de 3 letras y 3 dígitos, con repetición
        \item Números de 4 cifras distintas con los dígitos 1--9
        \item Ordenar 6 personas en fila
        \item Comités de 3 entre 6, sin cargos
      \end{enumerate}
    \end{column}
    \begin{column}{0.50\textwidth}
      \begin{enumerate}\setlength{\itemsep}{0.16cm}
        \setcounter{enumi}{10}
        \item Palabras distintas con BANANA
        \item Palabras distintas con ESTADISTICA
        \item Subconjuntos de un conjunto de 5 elementos
        \item Respuestas a 10 preguntas falso/verdadero
        \item Comidas: 3 entradas, 4 platos, 2 postres
        \item Equipos de 5 entre 12 jugadores
        \item Comisiones de 2 hombres (de 5) y 3 mujeres (de 6)
        \item Secuencias al lanzar una moneda 5 veces
        \item De las del 18, ¿cuántas con 3 caras exactas?
        \item $P(3$ caras exactas en 5 lanzamientos$)$
      \end{enumerate}
    \end{column}
  \end{columns}
\end{frame}

% ----------------------------------------------------------- RESPUESTAS ----
\begin{frame}{Respuestas}{Autoevaluación de los ejercicios propuestos}
  \listatiny
  \vspace{0.15cm}
  \begin{columns}[T,onlytextwidth]
    \begin{column}{0.47\textwidth}
      \begin{enumerate}\setlength{\itemsep}{0.16cm}
        \item $120$
        \item $336$ \;{\color{uniGris}$=P(8,3)$}
        \item $42$
        \item $21$
        \item $210$
        \item $210$ \;{\color{uniTitulo}igual al 5: $\binom{n}{r}=\binom{n}{n-r}$}
        \item $26^{3}\cdot 10^{3}=17\,576\,000$
        \item $P(9,4)=3024$
        \item $6!=720$
        \item $C(6,3)=20$
      \end{enumerate}
    \end{column}
    \begin{column}{0.50\textwidth}
      \begin{enumerate}\setlength{\itemsep}{0.16cm}
        \setcounter{enumi}{10}
        \item $6!/(3!\,2!)=60$
        \item $11!/(2!)^{4}=2\,494\,800$
        \item $2^{5}=32$
        \item $2^{10}=1024$
        \item $3\cdot 4\cdot 2=24$
        \item $C(12,5)=792$
        \item $C(5,2)\cdot C(6,3)=10\cdot 20=200$
        \item $2^{5}=32$
        \item $C(5,3)=10$
        \item $10/32=0{,}3125$
      \end{enumerate}
    \end{column}
  \end{columns}
\end{frame}

% ---------------------------------------------------------- REFERENCIAS ----
\begin{frame}{Referencias}
  \footnotesize
  \begin{itemize}
    \item[{[1]}] Díaz Rodríguez, M. (2022). \textit{Estadística inferencial
      aplicada}. Universidad del Norte.
      \texttt{https://elibro.net/es/lc/uniminuto/titulos/221614}
    \item[{[2]}] Llinás Solano, H. (2014). \textit{Introducción a la
      estadística con aplicaciones en ciencias sociales}. Universidad del
      Norte. \texttt{https://elibro.net/es/lc/uniminuto/titulos/70062}
    \item[{[3]}] Contento Rubio, M. R. (2019). \textit{Estadística con
      aplicaciones en R}. Editorial Utadeo.
      \texttt{https://elibro.net/es/lc/uniminuto/titulos/220926}
  \end{itemize}

  \vspace{0.35cm}
  \begin{center}
    {\Large\color{uniAzul}\textbf{¿Preguntas?}}
  \end{center}
\end{frame}

\end{document}
